% ═══════════════════════════════════════════════════════════════
% LEHRERHINWEISE: Kernreaktionen (Physik Klasse 10)
% Stunde 7-8 (90 Minuten)
% ═══════════════════════════════════════════════════════════════

\documentclass[11pt, a4paper]{article}

\usepackage[utf8]{inputenc}
\usepackage[T1]{fontenc}
\usepackage[ngerman]{babel}
\usepackage{helvet}
\renewcommand{\familydefault}{\sfdefault}

\usepackage[margin=2cm]{geometry}
\usepackage{enumitem}
\usepackage{tabularx}
\usepackage{booktabs}
\usepackage{array}
\usepackage{xcolor}
\usepackage{tcolorbox}
\usepackage{fancyhdr}

\setlist{nosep, leftmargin=1.5em}

% Farben
\definecolor{physik}{HTML}{2c5aa0}
\definecolor{grau}{HTML}{888888}
\definecolor{hellgrau}{HTML}{f0f0f0}
\definecolor{warnung}{HTML}{b35c00}
\definecolor{warnungbg}{HTML}{fff3e0}

% Kopfzeile
\pagestyle{fancy}
\fancyhf{}
\renewcommand{\headrulewidth}{0.4pt}
\lhead{\small\textbf{LEHRERHINWEISE} -- Physik Klasse 10}
\rhead{\small Kernreaktionen (Stunde 7-8)}
\cfoot{\small\thepage}

% Boxen
\newtcolorbox{warnbox}[1][]{
    colback=warnungbg, colframe=warnung, boxrule=0.8pt, arc=0pt,
    left=6pt, right=6pt, top=6pt, bottom=6pt,
    fonttitle=\bfseries\small, title=#1
}

\newcommand{\abschnitt}[1]{%
    \vspace{10pt}\noindent\textbf{\textcolor{physik}{\large #1}}\vspace{6pt}\hrule\vspace{8pt}
}

\begin{document}

\begin{center}
{\LARGE\bfseries Lehrerhinweise: Kernreaktionen}\\[6pt]
{\large Physik Klasse 10 -- Stunde 7-8 (90 Min)}\\[4pt]
{\small\textcolor{grau}{Kernphysik-Einheit, Teil 4}}
\end{center}
\vspace{8pt}\hrule\vspace{12pt}

% ═══════════════════════════════════════════════════════════════
\abschnitt{1. Stundenverlauf (90 Minuten)}

\begin{tabularx}{\textwidth}{|c|c|X|l|}
\hline
\textbf{Zeit} & \textbf{Phase} & \textbf{Inhalt} & \textbf{Material} \\
\hline
0-5 & Einstieg & Frage: Woher kommt die Energie der Sonne? & Tafel \\
\hline
5-20 & Input & Lernpfad: Kernspaltung (U-235) & LP, Beamer \\
\hline
20-30 & Erarbeitung & AB Kasten 1: Spaltungsgleichung & AB-04 \\
\hline
30-40 & Simulation & PhET Nuclear Fission: Kettenreaktion & Tablets/PCs \\
\hline
40-50 & Erarbeitung & AB Kasten 2 + Seite 2 (Protokoll) & AB-04 \\
\hline
50-60 & Input & Lernpfad: Kernfusion (Sonne, ITER) & LP \\
\hline
60-70 & Erarbeitung & AB Kasten 3, Übung 1-2 & AB-04 \\
\hline
70-80 & Diskussion & Kernkraftwerk vs. Atombombe (Übung 3) & Plenum \\
\hline
80-90 & Vertiefung & Übung 4: Vor-/Nachteile, Zusammenfassung & AB-04, Tafel \\
\hline
\end{tabularx}

\vspace{12pt}

% ═══════════════════════════════════════════════════════════════
\abschnitt{2. Differenzierung}

\begin{tabularx}{\textwidth}{|c|X|X|}
\hline
\textbf{Niveau} & \textbf{Fokus} & \textbf{Hilfestellung} \\
\hline
★ (BR) & Spaltung vs. Fusion zuordnen, qualitativ & Bildkarten Sonne/KKW \\
\hline
★★ (MR) & + Reaktionsgleichungen, Kettenreaktion & Formelkarte, Simulation \\
\hline
★★★ (GY) & + Kritische Masse, Energievergleich, Bewertung & Recherche zu ITER \\
\hline
\end{tabularx}

\vspace{12pt}

% ═══════════════════════════════════════════════════════════════
\abschnitt{3. Häufige Fehler \& Reaktionen}

\begin{warnbox}[Typische Schülerfehler]
\begin{enumerate}
\item \textbf{„Ein Kernkraftwerk kann wie eine Atombombe explodieren"}\\
→ Korrektur: Anreicherung im KKW nur 3-5\% (Bombe >90\%). Kontrollstäbe verhindern unkontrollierte Reaktion. Physikalisch unmöglich!

\item \textbf{Verwechslung Spaltung/Fusion}\\
→ Eselsbrücke: \textbf{Spaltung} = \textbf{S}chwere Kerne \textbf{s}palten. \textbf{Fusion} = \textbf{F}euer der Sonne (leichte Kerne).

\item \textbf{„Fusion ist fertig entwickelt"}\\
→ Klarstellung: Fusion auf der Erde noch nicht beherrscht (ITER = Experiment). Problem: 100 Mio. °C erzeugen und einschließen.

\item \textbf{Energiefreisetzung nicht verstanden}\\
→ E=mc²: Massendefekt wird zu Energie. Bindungsenergie-pro-Nukleon-Kurve zeigen.
\end{enumerate}
\end{warnbox}

\vspace{12pt}

% ═══════════════════════════════════════════════════════════════
\abschnitt{4. Benötigtes Material}

\begin{itemize}
\item Lernpfad: \texttt{lernpfad-kernreaktionen.html}
\item Arbeitsblatt: \texttt{AB-04-kernreaktionen.pdf} (Klassensatz)
\item Musterlösung: \texttt{ML-04-kernreaktionen.pdf}
\item PhET-Simulation: \texttt{phet.colorado.edu/de/simulations/nuclear-fission}
\item Optional: Video „Wie funktioniert ein Kernkraftwerk?" (SimpleClub o.ä.)
\item Optional: Bilder ITER-Projekt, Tokamak
\end{itemize}

\vspace{12pt}

% ═══════════════════════════════════════════════════════════════
\abschnitt{5. Lösungen (Kurzfassung)}

\textbf{Kasten 1:} Kernspaltung durch Neutron. U-235 + n → Ba-141 + Kr-92 + 3n + Energie

\textbf{Kasten 2:} Kettenreaktion: Neutronen spalten weitere Kerne. Kontrolliert = KKW, Unkontrolliert = Bombe. Kritische Masse = Mindestmenge für selbsterhaltende Reaktion.

\textbf{Kasten 3:} Fusion: Leichte Kerne verschmelzen. D + T → He-4 + n + Energie. Benötigt 100 Mio. °C.

\textbf{Übung 1:} Sonne=F, U-235=S, Verschmelzung=F, KKW=S, 100 Mio. °C=F, Radioaktiver Müll=S

\textbf{Übung 2a:} $^{235}_{92}$U + n → $^{92}_{36}$Kr + $^{141}_{56}$Ba + 3n

\textbf{Übung 2b:} 1→3→9→27→81 Neutronen (oder $3^4$)

\vspace{8pt}
\textit{→ Vollständige Lösungen siehe ML-04-kernreaktionen.pdf}

\end{document}
